\documentclass[10pt,letterpaper,twocolumn]{article}

\usepackage[T1]{fontenc}
\usepackage{lmodern}
\usepackage[margin=0.64in,columnsep=0.27in]{geometry}
\usepackage{amsmath,amssymb,mathtools,bm}
\usepackage{array,booktabs,tabularx}
\usepackage{xcolor}
\usepackage{enumitem}
\usepackage{microtype}
\usepackage{fancyhdr}
\usepackage{hyperref}
\usepackage{needspace}

% The user-local tools/array package is newer than the available LaTeX kernel.
\makeatletter
\def\vcenter@text{\vbox}
\makeatother

\definecolor{navy}{HTML}{17365D}
\definecolor{muted}{HTML}{4A5568}
\hypersetup{
  colorlinks=true,
  linkcolor=navy,
  urlcolor=navy,
  pdftitle={M4 Moment-Hierarchy Solutions and Debrief},
  pdfauthor={Local Paper-V learning sequence}
}

\pagestyle{fancy}
\fancyhf{}
\fancyhead[L]{\footnotesize $M_4$ moment hierarchy}
\fancyhead[R]{\footnotesize Solutions and debrief}
\fancyfoot[C]{\thepage}
\renewcommand{\headrulewidth}{0.3pt}

\setlength{\parindent}{0pt}
\setlength{\parskip}{3pt plus 0.5pt}
\setlist{nosep,leftmargin=*,topsep=2pt}
\renewcommand{\arraystretch}{1.12}
\newcolumntype{Y}{>{\raggedright\arraybackslash}X}
\newcommand{\R}{\mathbb R}
\newcommand{\C}{\mathbb C}
\newcommand{\ii}{\mathrm i}
\newcommand{\Tr}{\operatorname{Tr}}
\newcommand{\expect}[1]{\left\langle #1\right\rangle}
\newcommand{\sectionline}[1]{%
  \Needspace{6\baselineskip}%
  \par\medskip{\color{navy}\large\bfseries #1}\par
  \vspace{-2pt}{\color{navy}\rule{\columnwidth}{0.7pt}}\par\smallskip
}

\begin{document}
\raggedbottom

\twocolumn[
\begin{center}
{\color{navy}\LARGE\bfseries From the Archive Equations to the $M_4$ Cone}\par
\vspace{3pt}{\Large Solutions, Reference Reconstructions, and Debrief Criteria}\par
\vspace{4pt}{\small Companion to the closed-book workbook, 5 August 2026}
\end{center}
\vspace{3pt}\hrule\vspace{7pt}
]

\textbf{Use condition.}  Preserve the first attempt and make a second attempt
before reading the corresponding solution.  A displayed solution is a
reference reconstruction, not evidence that the construction can yet be
reproduced.  The transfer-and-repair problems remain intentionally unsolved.

\sectionline{Part I --- Reference reconstructions}

\textbf{R1. Governing map.}
The initial reduction is
\[
(H,|\Psi_0\rangle)\longrightarrow X_0
\longrightarrow u^0=x_{{\rm raw},0}\oplus\eta_0\in\R^{60}.
\]
At each right-hand-side evaluation,
\[
u\xrightarrow{\mathcal C_4}m_4^\star
\longrightarrow f_{\le3}
\xrightarrow{K/P/D}f_{\rm aug}
\xrightarrow{\mathcal C_{\rm ret}}f_{\rm corrected}.
\]
Here \(\mathcal C_4\) is degree-four positive completion and
\(\mathcal C_{\rm ret}\) is the optional retained velocity slice.  The time
step then follows
\[
f_{\rm corrected}\xrightarrow{\mathrm{SSPRK}(3,3)}
\widetilde u^{(j)}\xrightarrow{R_{M_4}}u^{(j)},
\]
where \(R_{M_4}\) is the full spectral check and hidden-stage retraction.  At
the end,
\[
u(t)\longrightarrow\mathcal O(t)
\longrightarrow\text{time-RMS errors}.
\]
The exact trajectory is not an online controller input.  Its contractions
prepare the common initial reduced state, and its later trajectory is used
offline for observable scoring.  The RMS contraction and comparison are also
offline.

\textbf{R2. Expectation-value equation.}
From \(\ii|\dot\Psi\rangle=H|\Psi\rangle\),
\(|\dot\Psi\rangle=-\ii H|\Psi\rangle\) and
\(\langle\dot\Psi|=\ii\langle\Psi|H\).  Therefore
\[
\begin{aligned}
\frac{d}{dt}\langle O\rangle
&=\langle\dot\Psi|O|\Psi\rangle
 +\langle\Psi|O|\dot\Psi\rangle\\
&=\ii\langle HO\rangle-\ii\langle OH\rangle
=\ii\langle[H,O]\rangle.
\end{aligned}
\]
For \(O(t)\), add \(\langle\partial_tO\rangle\).

\textbf{R3. One upward chain.}
With
\[
H_-=-t(X_\uparrow+X_\downarrow)
+\frac{\Delta}{2}(Z_\uparrow+Z_\downarrow)
+\frac{\omega}{2}(x^2+p^2)+gx(Z_\uparrow+Z_\downarrow),
\]
Pauli commutators give
\[
\dot{\expect{Y_s}}
=2t\expect{Z_s}+\Delta\expect{X_s}
+2g\expect{xX_s}.
\]
The degree-one moment \(\expect{Y_s}\) requests the degree-two moment
\(\expect{xX_s}\).  Its derivative is
\[
\frac{d}{dt}\expect{xX_s}
=\omega\expect{pX_s}-\Delta\expect{xY_s}
-2g\expect{x^2Y_s},
\]
which requests the degree-three moment \(\expect{x^2Y_s}\).  The interaction
term in its next derivative can request a degree-four moment such as
\(\expect{x^3X_s}\).  Thus the chain is generated by the Hamiltonian algebra,
not by the integrator.

\textbf{R4. Coordinate counts.}
The archive state contributes \(\rho:3\), \(B:4\), \(N:4\), \(A:6\), and
\(C:14\), so \(3+4+4+6+14=31\).
Here \(\rho\) is Hermitian with unit trace; \(B\) has two complex entries;
\(N\) is Hermitian; \(A\) is complex symmetric; and the two complex
\(2\times2\) correlation slices obey
\(\Tr C^0=\Tr C^1\), reducing their unrestricted 16 real entries to 14.
The APCM entrance layer imposes the stronger centered condition
\(\Tr C^0=\Tr C^1=0\).  Setting the remaining shared complex trace to zero
removes two real coordinates, hence \(31-2=29\).  Six two-spin electronic
moments plus all 25 degree-three moments form the 31 hidden coordinates, so
\(29+31=60\).

\textbf{R5. Moment positivity.}
For \(Q(z)=\sum_a z_aW_a\), define
\(M_{ab}=\expect{W_a^\dagger W_b}\).  Then
\[
z^\dagger Mz
=\sum_{ab}z_a^*\expect{W_a^\dagger W_b}z_b
=\expect{Q(z)^\dagger Q(z)}\ge0.
\]
The last inequality follows because \(Q^\dagger Q\) is positive.  It holds
for every \(z\), so \(M\succeq0\).  If each word has degree at most two,
the product \(W_a^\dagger W_b\) can have degree four; hence its Gram entry can
require fourth moments.

\textbf{R6. Three distinct operations.}
\begin{center}
\footnotesize
\begin{tabularx}{\columnwidth}{@{}lYYY@{}}
\toprule
operation & variable & PSD object & stage seam\\
\midrule
retained slice & velocity \(w\in\R^{31}\) & \(\rho, I-\rho,\mathcal G\) barrier derivatives & inside each RHS evaluation\\
completion & frontier \(m_4\in\R^{35}\) & \(M_4(\mu_{\le3},m_4)\) & before the degree-three RHS\\
retraction & hidden \(\eta'\in\R^{31}\), with a witness \(m_4\) & the same \(M_4\) & after every provisional SSPRK stage\\
\bottomrule
\end{tabularx}
\end{center}
Completion supplies a velocity input; retraction modifies a provisional lower
state; retained slicing modifies the archive velocity.  They are not three
names for one projection.

\sectionline{Part II --- Exit-quiz key}

\begin{center}
\begin{tabular}{@{}lcccccccc@{}}
\toprule
item & C1&C2&C3&C4&C5&C6&C7&C8\\
\midrule
answer & C&A&D&B&C&A&D&B\\
\bottomrule
\end{tabular}
\end{center}
\begin{center}
\begin{tabular}{@{}lcccc@{}}
\toprule
item & T1&T2&T3&T4\\
\midrule
answer & C&A&B&D\\
\bottomrule
\end{tabular}
\qquad
\begin{tabular}{@{}lccccc@{}}
\toprule
item & M1&M2&M3&M4&M5\\
\midrule
answer & B&A&D&C&B\\
\bottomrule
\end{tabular}
\end{center}

\textbf{C1--C4.}  C1 is C because a retained derivative requests data outside
the retained state; closure supplies that data.  C2 is A because vanishing a
terminal cumulant is a deterministic modeling rule, not a theorem about the
exact state or positivity.  C3 is D because the archive backbone and hidden
moments carry complementary coordinates.  C4 is B because \(M_4\) is the
word Gram matrix, not a fourth-cumulant vector.

\textbf{C5--C8.}  C5 is C: completed fourth moments feed the third-degree
velocity, and retraction can alter hidden stages.  C6 is A: the retained
controller finds a nearby equality-compatible inward velocity.  C7 is D:
the observation is local to the matched protocol and horizon.  C8 is B:
exact propagation supplies common initialization and offline scoring, not
online corrections.

\textbf{T1--T4.}  A spectrahedron is an affine PSD slice (T1 C).  Retraction
is a scaled nearest change of selected hidden lower moments that restores an
extension (T2 A).  The scalar \(v^\dagger\Delta B(e_j)v\) is a first-order
directional response (T3 B).  A \(10^{-6}\) conic backward-error tolerance
licenses only eigenvalue-level numerical feasibility at that scale (T4 D).

\textbf{M1--M5.}  M1 is the Heisenberg identity (B).  M2 has two nonidentity
Pauli factors and boson degree three, hence total degree five (A).  M3 is the
positive-square identity (D).  M4 is the orthogonal displacement to the
supporting hyperplane (C).  M5 is the exact affine cut for the exposed fixed
mode (B).

\sectionline{Part III --- Derivations}

\textbf{D1. Zero third cumulant.}
For commuting placeholders,
\[
\begin{aligned}
\kappa(A,B,C)={}&\expect{ABC}
-\expect{AB}\expect C-\expect{AC}\expect B
-\expect{BC}\expect A\\
&+2\expect A\expect B\expect C.
\end{aligned}
\]
Setting \(\kappa(A,B,C)=0\) gives
\[
\expect{ABC}=\expect{AB}\expect C+\expect{AC}\expect B
+\expect{BC}\expect A-2\expect A\expect B\expect C.
\]
At a degree-three truncation, the analogous partition rule sets each omitted
degree-four connected cumulant to zero and reconstructs its raw fourth moment
from retained lower moments.  In the implementation this produces the prior
\(m_4^{(0)}\), which positive completion may move.

\textbf{D2. Relative oscillator.}
Only the oscillator and coupling terms contribute.  Using
\([p^2,x]=-2\ii p\),
\[
\dot{\expect x}
=\ii\expect{[\tfrac\omega2p^2,x]}
=\ii\frac\omega2(-2\ii)\expect p
=\omega\expect p.
\]
Using \([x^2,p]=2\ii x\) and \([xZ_s,p]=\ii Z_s\),
\[
\begin{aligned}
\dot{\expect p}
&=\ii\expect{[\tfrac\omega2x^2+gx\sum_sZ_s,p]}\\
&=-\omega\expect x-g\expect{Z_\uparrow+Z_\downarrow}.
\end{aligned}
\]

\textbf{D3. One-cut KKT solution.}
Write the constraint as \(g(w)=-(a^{\mathsf T}w+\lambda)\le0\) and
\[
\mathcal L=\tfrac12w^{\mathsf T}w-\mu(a^{\mathsf T}w+\lambda),
\qquad \mu\ge0.
\]
Stationarity gives \(w=\mu a\); feasibility and complementarity give
\(a^{\mathsf T}w+\lambda\ge0\) and
\(\mu(a^{\mathsf T}w+\lambda)=0\).  If \(\lambda\ge0\), \(w=0\) is feasible
and optimal, with \(\mu=0\).  If \(\lambda<0\), the constraint is active:
\[
\mu=-\frac{\lambda}{a^{\mathsf T}a},\qquad
w^\star=-\frac{\lambda}{a^{\mathsf T}a}a.
\]
The multiplier is the distance scale required by the violated mode; \(a\) is
the correction-space normal that most efficiently raises that fixed quadratic
form.  If \(a=0\) while \(\lambda<0\), this correction chart cannot repair the
mode.

\textbf{D4. Equalities first.}
Let the compact SVD be \(L=U_r\Sigma_rV_r^{\mathsf T}\).  For consistent
equalities, the minimum-norm particular solution is
\[
w_{\rm p}=L^+d=V_r\Sigma_r^{-1}U_r^{\mathsf T}d.
\]
Let the columns of \(Z\) be an orthonormal basis for \(\ker L\).  Then
\(LZ=0\), \(Z^{\mathsf T}Z=I\), and every solution is
\(w=w_{\rm p}+Zy\).  Since the minimum-norm particular solution lies in
\(\operatorname{range}(L^{\mathsf T})\), it is orthogonal to \(\ker L\), so
\[
\frac12\|w\|_2^2
=\frac12\|w_{\rm p}\|_2^2+\frac12\|y\|_2^2.
\]
Thus the equality cost is constant and the cone problem is Euclidean in \(y\).

\textbf{D5. One slice versus the cone.}
For a fixed normalized \(v\) and affine
\(B(y)=B_0+\sum_jy_jB_j\),
\[
\begin{aligned}
v^\dagger B(y)v
&=v^\dagger B(y^{(k)})v\\
&\quad+\sum_j(y_j-y_j^{(k)})v^\dagger B_jv\\
&=\lambda+n^{\mathsf T}(y-y^{(k)}).
\end{aligned}
\]
The cut \(n^{\mathsf T}y\ge n^{\mathsf T}y^{(k)}-\lambda\) is therefore exact
for that fixed direction.  But
\[
\lambda_{\min}(B(y))=\min_{\|q\|=1}q^\dagger B(y)q.
\]
Satisfying the cut controls only \(q=v\); another direction can become the
minimizer and remain negative.  Certification requires rebuilding every full
barrier matrix and checking that every minimum eigenvalue is no smaller than
the declared negative tolerance.

\textbf{D6. Factorial ablation.}
Relative to the no-cone higher-moment cell, the $M_4$-only improvement factors
are
\[
\begin{aligned}
0.04995/0.01258&=3.97,\\
0.10663/0.04743&=2.25,\\
0.02509/0.000821&=30.56.
\end{aligned}
\]
Mechanistic conclusion: positive completion is dynamically active in this
protocol because it supplies fourth moments used by the velocity and supports
stage retractions; the large error changes are consistent with that mechanism.
Unsupported conclusion: the table does not show an exact closure, a universal
implication from \(M_4\) positivity to accuracy, transfer to other parameters,
or behavior beyond \(t t_{\rm hop}=4\).

\textbf{D7. SSPRK placement.}
For \(F(t,u)\) denoting a completed and, when enabled, retained-corrected RHS,
\[
\begin{aligned}
\widetilde u^{(1)}&=u^n+hF(t_n,u^n),\\
u^{(1)}&=R_{M_4}(\widetilde u^{(1)}),\\
\widetilde u^{(2)}&=\tfrac34u^n+\tfrac14
 [u^{(1)}+hF(t_n+h,u^{(1)})],\\
u^{(2)}&=R_{M_4}(\widetilde u^{(2)}),\\
\widetilde u^{n+1}&=\tfrac13u^n+\tfrac23
 [u^{(2)}+hF(t_n+\tfrac12h,u^{(2)})],\\
u^{n+1}&=R_{M_4}(\widetilde u^{n+1}).
\end{aligned}
\]
Each evaluation of \(F\) first completes \(m_4\), uses it in the degree-three
velocity, and then adds the retained controller's archive-velocity correction
when that switch is on.  The map \(R_{M_4}\) performs a full spectral check and
hidden-state retraction only when needed.  End-of-step-only checking omits the
two intermediate retractions; those uncorrected stage states change later RHS
evaluations, so it defines a different discrete dynamical map.

\sectionline{Part IV --- Debrief criteria only}

The transfer problems X1--X5 deliberately have no worked answers here.  Score
them by whether the response contains the following kinds of evidence, without
using this list to repair the first attempt.

\begin{enumerate}
\item \textbf{X1:} an explicit lost mixed operator direction, a valid
\(Q(z)\) comparison between the full and reduced word spans, and a bounded
claim about what the reduced Gram can no longer test.
\item \textbf{X2:} a concrete positive yet inaccurate reduced sequence and a
clean separation between representability, dynamical accuracy, and numerical
stability.
\item \textbf{X3:} a positive-definite metric, its square-root or Cholesky
factor, the transformed equality/cone maps, and an explanation of why metric
choice selects a different nearest feasible point.
\item \textbf{X4:} one explanation based on unequal word content and another
based on inconsistent lifts or coordinates, with a diagnostic that can
distinguish them.
\item \textbf{X5:} pointwise completion and retraction parity, trajectory
parity, event/tolerance parity, observable parity, and a reason near-degenerate
optimizer selection can matter even under an identical feasible set.
\end{enumerate}

Do not convert a failed transfer item into ``mastered'' by reading nearby
source text.  Record the earliest unsupported construction, revisit only that
dependency, and retry the original item with its wording unchanged.

\sectionline{Scoring record}

\begin{center}
\begin{tabularx}{\columnwidth}{@{}Yccc@{}}
\toprule
evidence class & independent & after hint & unresolved\\
\midrule
governing-map reconstruction &&&\\[8pt]
concept and terminology &&&\\[8pt]
local derivations &&&\\[8pt]
transfer and repair &&&\\[8pt]
\bottomrule
\end{tabularx}
\end{center}

Mastery of this route requires independent reconstruction of the hierarchy and
cone algorithms plus successful transfer to changed word bases, metrics, or
solver semantics.  Completion of the reading alone records exposure.

\end{document}
